% Section II: Topological structures already present in ordinary gauge theory
% Purpose: Bridge from Schwartz - continuation of standard QFT, not a jump into pure math
% Make reader feel: "I already know the raw materials. What changes now is how they are organized."
%
% Content checklist:
% - Yang-Mills vacua and large gauge transformations
% - π₃(G), instanton number
% - θ-term and topological sectors
% - Anomalies and descent equations
% - Wilson loops as nonlocal observables
% - Monopoles and Dirac charge quantization
%
% Status: NEEDS TO BE WRITTEN (some content may exist in class notes)
% Sources:
% - Review PROJECTS/QFT/literature/ for class materials on instantons/monopoles/θ-vacua
% - May reference Schwartz chapters as needed
%
% End goal statement:
% "Ordinary QFT contains topological sectors, but still has local propagating
%  degrees of freedom and metric dependence. We therefore look for a theory
%  in which the topological data are themselves the primary degrees of freedom."
%
% Writing notes:
% - This is NOT a full review - just enough to establish what student already knows
% - Each subsection: quick reminder + "this is topological structure" + "but theory is not TQFT"
% - ~10-15 pages total

\section{Topological Structure in Ordinary Gauge Theory}

% TODO: Write new content reviewing standard QFT topological features
% Keep it survey-level but concrete
% Each example should end with "but the theory still has local dynamics and metric dependence"

\subsection{Large Gauge Transformations and Vacuum Structure}

% TODO: π₃(G), instanton sectors, vacuum angle

\subsection{\texorpdfstring{The $\theta$-Term}{The theta-Term}}

% TODO: ∫F∧F̃, total derivative, topological but not TQFT

\subsection{Anomalies and Descent}

% TODO: Preview anomaly descent without full derivation (that comes in Section IV)

\subsection{Wilson Loops}

% TODO: Nonlocal observables, but still in dynamical gauge theory

\subsection{Magnetic Monopoles}

% TODO: Dirac quantization, topological charge, but solitons in ordinary QFT

\subsection{Transition to Topological Quantum Field Theory}

% TODO: Summary pointing toward CS theory
% "We now ask: what if topology itself controls the physically meaningful data?"
